% Generated draft: requires mathematical and source review.
\documentclass[11pt]{article}
\usepackage[margin=25mm]{geometry}
\usepackage{fontspec}
\setmainfont{Latin Modern Roman}
\usepackage{amsmath,amssymb,mathtools,unicode-math}
\setmathfont{Latin Modern Math}
\usepackage{enumitem,needspace,xurl,hyperref}
\setlength{\parindent}{0pt}
\setlength{\parskip}{5pt}
\setlength{\emergencystretch}{3em}
\newcommand{\N}{\mathbb{N}}
\newcommand{\Z}{\mathbb{Z}}
\newcommand{\Q}{\mathbb{Q}}
\newcommand{\R}{\mathbb{R}}
\newcommand{\C}{\mathbb{C}}
\newcommand{\F}{\mathbb{F}}
\newcommand{\A}{\mathbb{A}}
\newcommand{\PP}{\mathbb{P}}
\newcommand{\E}{\mathbb{E}}
\newcommand{\eps}{\varepsilon}
\newcommand{\dd}{\,\mathrm{d}}
\begin{document}
\section*{UnsolvedMath ID: 30006193}
\textbf{Source problem code:} catalogue entry
\subsection{Definitions and mathematical statement}
Let \(F\) be a local field, Archimedean or non-Archimedean, and let \(\mathbf G\) be a semisimple algebraic group defined over \(F\). Put \(G=\mathbf G(F)\), with its local-field topology, and let
\[
r=\operatorname{rk}_F(\mathbf G),
\]
the dimension of a maximal \(F\)-split torus of \(\mathbf G\). No positive-rank assumption is stated, so \(r=0\) is included.

For \(1<p<\infty\), let \(L^p(G)\) denote the space of Haar-measurable functions modulo equality almost everywhere whose absolute value has integrable \(p\)-th power, with respect to a left Haar measure on \(G\). To make the notation explicit, use the conventional left regular action
\[
(g\cdot f)(x)=f(g^{-1}x).
\]
Here the scalar field may be \(\mathbb R\) or \(\mathbb C\); the supplied source does not specify it. Interpret \(H^k_{\mathrm{ct}}(G,V)\) as ordinary continuous group cohomology: the cohomology of the inhomogeneous complex \(C^k_{\mathrm{ct}}(G,V)=C(G^k,V)\), with \(C^0_{\mathrm{ct}}(G,V)=V\), using the differential
\[
\begin{aligned}
(dc)(g_1,\ldots,g_{k+1})
={}&g_1\cdot c(g_2,\ldots,g_{k+1})\\
&+\sum_{i=1}^{k}(-1)^i c(g_1,\ldots,g_i g_{i+1},\ldots,g_{k+1})\\
&+(-1)^{k+1}c(g_1,\ldots,g_k).
\end{aligned}
\]
Thus \(H^k_{\mathrm{ct}}=\ker d/\operatorname{im}d\), without taking the closure of the coboundaries, and \(H^0_{\mathrm{ct}}(G,V)=V^G\).

For every such \(F\) and \(\mathbf G\), the supplied problem asks both:
\begin{enumerate}
\item Is it true that
\[
\forall p\in(1,\infty)\;\forall k\in\mathbb Z_{\geq 0},\qquad
k<r\ \Longrightarrow\ H^k_{\mathrm{ct}}(G,L^p(G))=\{0\}?
\]
\item Is it true that
\[
\exists p\in(1,\infty)\qquad H^r_{\mathrm{ct}}(G,L^p(G))\neq\{0\}?
\]
The exponent in the second question may depend on \(F\) and \(\mathbf G\).
\end{enumerate}
When \(r=0\), the first question has no eligible cohomological degrees, and the second concerns invariant vectors in degree zero.

The coefficient scalar field, regular-action convention, and distinction between ordinary and reduced cohomology are not explicitly defined in the supplied passage; the conventions above make this draft readable and require source confirmation. The source also does not spell out whether connectedness is part of its convention for a semisimple algebraic group. No additional assumptions such as simplicity, isotropy, or simple connectedness have been inserted.
\subsection{Short English statement}
For the group of points of a semisimple algebraic group over any local field, ask whether continuous cohomology with regular \(L^p\) coefficients vanishes in every nonnegative degree below the \(F\)-rank for every \(1<p<\infty\), and whether it is nonzero in the rank degree for at least one such \(p\). The supplied record labels the problem partially solved; coefficient and cohomology conventions require confirmation.
\subsection{Sources}
\begin{itemize}
\item Local UnsolvedMath registry, supplied record with permanent ID \(30006193\), title \(L^p\) Cohomology at the Semisimple Rank Threshold, source problem code \texttt{catalogue entry}, in the Oberwolfach Reports Open Problems collection. The supplied original excerpt begins with the marker \texttt{253]}. No individual proposer is identified; the supplied proposed year is 2025.
\item Source report: \emph{Geometric Structures in Group Theory} (2025), Oberwolfach Reports, pp.\ 275--338, DOI \texttt{10.4171/owr/2025/6}. \url{https://doi.org/10.4171/owr/2025/6}. The registry labels the statement \texttt{corrected\_verified} and says it was reconstructed and checked against the official report. This is supplied verification metadata, not an independent verification performed for this draft.
\item Supplied literature reference: Bourdon and Remy, nonvanishing \(L^p\)-cohomology. \url{https://doi.org/10.5802/jep.232}.
\item Supplied literature reference: Lopez Neumann, vanishing in degree two. \url{https://arxiv.org/abs/2302.09307}.
\item Supplied literature reference: OWR 2025/6 geometric group theory report. \url{https://oa.tib.eu/renate/server/api/core/bitstreams/6f61d540-8e50-41f3-97a8-6e8214672460/content}.
\item Supplied literature assessment, checked \textbf{2026-08-22}: status \emph{partially solved}. The supplied assessment states that, for real and non-Archimedean semisimple groups, substantial vanishing below the rank and nonvanishing in the rank degree are proved in important ranges, including large \(p\) and many group types, while the all-local-field, all-\(p\) formulation remains open. This is supplied dated curation derived from the OWR record, not an independently verified literature assessment. The supplied background expressly calls for independent verification of time-sensitive status and cited claims.
\end{itemize}
\end{document}
